\section{Related Work}

\noindent\textbf{Mamba for long sequence modeling.}~~
The structured State-Space Sequence (S4) model is a novel alternative to CNNs or transformers to model the long-term dependency \cite{12}. The promising property of linearly scaling in sequence length attracts further exploration. Based on the S4, \citet{16} propose a new S5 layer by introducing MIMO SSM and an efficient parallel scan into the S4 layer. \citet{17} design a new SSM layer, H3, that nearly fills the performance gap between SSMs and transformers in language modeling. \citet{18} build the Gated State Space layer on S4 by introducing more gating units to improve the expressivity.
Recently, \citet{10} propose a data-dependent SSM layer and build a generic language model backbone, Mamba, which outperforms transformers at various sizes on large-scale real data and enjoys linear scaling in sequence length. Later, Vision Mamba adds position coding and bidirectional scanning to extend it to visual tasks \cite{11}. 
In this work, we explore transferring Mamba's success to RL, i.e., achieving high effectiveness and efficiency for long-term memory.

\noindent\textbf{Transformer for Decision-Making.}~~
In general, reinforcement learning was proposed as a fundamental online paradigm \cite{38}. The nature of online learning comes with some limitations when meeting the applications for which it is impossible to gather online data and learn simultaneously, such as autonomous driving. To this end, offline RL proposed that the agent can learn from a fixed offline dataset without gathering new data during learning \cite{39,40,41,42}. In the context of offline RL, recent works explored using transformer-based policy by treating RL tasks as a type of sequential prediction problem. Among them, a decision transformer \cite{2} was proposed to model trajectories as sequences and autoregressively predict action conditioning on desired return-to-go, past states, and actions. Trajectory transformer \cite{43} demonstrated that transformer could learn single-task policies from offline data. Subsequently, the multi-game decision transformer \cite{14} and Gato \cite{27} further showed that transformer-based policies could address multi-tasks in the same domain and cross-domain tasks. However, these works focused on distilling expert policies from offline data and failed to enable self-improvement like DM-H. When the offline data are sub-optimal, or the agent is required to adapt to new tasks, the multi-game decision transformers need to finetune the model parameters, while Gato is required to get prompted with expert demonstrations.

\noindent\textbf{Meta RL.}~~
DM-H falls into the category of methods of learning to learn, which is also known as meta-learning. More precisely, recent in-context RL methods can be categorized as in-context meta-RL methods. The general idea of learning self-improvement has a long history in RL but is limited to hyper-parameters in the early stages \cite{31}. In-context meta-RL methods \cite{33,32} are commonly trained in the online setting by maximizing multi-episodic value functions with memory-based architectures through environment interactions. Another online meta-RL attempts to find good network parameter initializations and then quickly adapt through additional gradient updates \cite{46,45}. More recently, meta-RL has seen substantial breakthroughs, from performance gains on popular benchmarks to offline settings, such as Bayesian RL \cite{36} and optimization-based meta-RL \cite{37}. Considering the difficulty of a completely offline setting, recent work has explored hybrid offline-online settings \cite{34,35}. DM-H is similar to the hybrid offline-online setting but saves more computing resources because the online phase does not involve gradient updates.

\noindent\textbf{In-Context RL.}~~
In-context RL is the one that addresses tasks by providing prompts or demonstrations \cite{2,43}. By training agents at a large scale, transformer-based policies usually have the ability to learn in context \cite{14,27}. The learning process is performed entirely in context and does not involve parameter updates of neural networks. In this work, we consider incremental in-context RL, which involves learning from one's own behaviors in a trial-and-error manner. \citet{8} proposed Algorithm Distillation (AD), which automatically improved its performance by providing multiple historical trajectories. Subsequently, \citet{3} proposed a Decision-Pretrained Transformer, which trained the agent to find optimal behaviors faster by only predicting the optimal trajectory. More recently, \citet{1} further demonstrated that across-episodic contexts encourage large transformer models' emerging self-improvement behaviors. However, these methods suffer from huge computational costs as across-episodic contexts induce too-long sequences. In contrast, DM-H leverages Mamba's ability to efficiently process long-term dependencies while using the transformer to establish high-quality predictions.

